% object_id: TEX-V22-P4-T02-B2-POST-CURRENT-SOURCE-KSTAR-THEOREM-ROUTE-SELECTION-V1
% authority: science_draft
% task_id: RT-20260814-003
\documentclass[11pt]{article}
\usepackage[margin=1in]{geometry}
\usepackage{amsmath,amssymb,booktabs,array,enumitem,microtype,url}
\setlist{nosep}
\newcommand{\Diff}{\operatorname{Diff}}
\newcommand{\Diffc}{\operatorname{Diff}_{c}}
\newcommand{\Kstar}{K_{\star}}

\title{V22 P4-T02/B2 Post-$K_{\star}$\\Theoretical Continuation Selection}
\author{The \AE ther-Flow Research Control Program}
\date{14 August 2026}

\begin{document}
\maketitle

\begin{abstract}
The exact predecessor theorem defines a diffeomorphism-natural wide normal
arrow-subclass assignment \(\Kstar\) on the ordinary smooth four-manifold
arena.  This note compares four bounded continuations and selects one future
source-side irrelevance theorem packet.  The selected packet tests whether the
unchanged assignment has restriction-compatible local bridge force.  It is
not executed here.  No equivalence law, physical allowedness, response
semantics, causal cone, effective metric, ontology adoption, or
Distance-to-GR advance is asserted.
\end{abstract}

\section{Fixed predecessor theorem and semantic boundary}

The fixed RT002 assignment is
\begin{equation}
\Kstar(M)=
\begin{cases}
\Diffc(M),&\Diffc(M)\subsetneq\Diff(M),\\
\{\operatorname{id}_M\},&\Diffc(M)=\Diff(M).
\end{cases}
\label{eq:kstar}
\end{equation}
For every diffeomorphism \(f:M\to N\), conjugation transports the selected
subgroup:
\begin{equation}
 f\Kstar(M)f^{-1}=\Kstar(N).
\end{equation}
This is a current-smooth-source mathematical theorem.  It is not a theorem
that \(\Kstar\) is \texttt{EqSrc}, physical allowedness, gauge, a detector
response map, a causal structure, or an effective-metric bridge.

\section{New exact decision controls}

\subsection{Compact/open localization mismatch}

Let \(M\) be a compact smooth four-manifold and let \(U\subset M\) be a
noncompact coordinate ball.  Compactness gives
\(\Diffc(M)=\Diff(M)\), so Equation~\eqref{eq:kstar} yields
\(\Kstar(M)=\{\operatorname{id}_M\}\).  On \(U\), compactly supported bump
diffeomorphisms are nonidentity and \(\Diffc(U)\) is proper in
\(\Diff(U)\), hence \(\Kstar(U)=\Diffc(U)\).

Choose nonidentity \(h\in\Diffc(U)\).  Extension by the identity gives a
smooth \(\widetilde h\in\Diff(M)\), but
\(\widetilde h\notin\Kstar(M)\).  Therefore the restrictions of global
\(\Kstar(M)\) arrows do not exhaust \(\Kstar(U)\).  Objectwise conjugation
naturality alone is not a restriction or gluing law.

This conclusion is deliberately scoped.  It does not show that every enriched
localization is impossible; it identifies a missing structure in the printed
objectwise assignment.

\subsection{Local first-jet isotropy}

Assume a chart stabilizer in \(\Diffc(U)\) realizes all positive-determinant
first jets \(GL^+(4)\) at a point.  If a tangent vector is invariant, the
positive scaling \(2I\) gives \(2v=v\), hence \(v=0\).  The same pullback
calculation forces any invariant covector to vanish.  Moreover,
\(GL^+(4)\) is transitive on nonzero tangent vectors.  Any invariant subset
of the punctured tangent space is therefore empty or the entire punctured
space, and cannot be a proper Lorentzian conformal cone.

The hypothesis and conclusion are representation-theoretic.  They neither
define physical time orientation nor identify a detector or propagation law.

\subsection{Large quotient versus local data}

Where \(\Diffc(M)\) is normal and the declared object class supports the
quotient, \(\Diff(M)/\Diffc(M)\) distinguishes large or end behavior.  It can
distinguish representatives whose germ on a fixed local cell agrees.  Thus a
large-diffeomorphism coset is not, without a separately defined natural map, a
local principal-symbol or response selector.

\subsection{Bridge relevance criterion}

A future claim that \(\Kstar\) has P4-T02 bridge force must provide all of:
\begin{enumerate}
  \item an independently current-source-defined local codomain;
  \item a typed natural local map into that codomain;
  \item restriction, overlap, composition, and gluing compatibility;
  \item noncollapse through the full local-diffeomorphism orbit quotient; and
  \item a claim boundary separating mathematics from physical and empirical
  interpretation.
\end{enumerate}
This is a necessary-condition audit, not a declaration that such a map cannot
exist.

\section{Four-route comparison}

\begin{center}
\small
\begin{tabular}{>{\raggedright\arraybackslash}p{0.11\linewidth}
                >{\raggedright\arraybackslash}p{0.23\linewidth}
                >{\raggedright\arraybackslash}p{0.23\linewidth}
                >{\raggedright\arraybackslash}p{0.31\linewidth}}
\toprule
Route & Packet class & Disposition & Decisive reason\\
\midrule
A & source-side irrelevance theorem & selected, not executed & Tests the first
missing locality and bridge-factor premises without adding source law data.\\
B & source-extension candidate & not selected & No independently sourced
P4-relevant local codomain or map is currently formed; construction risks
relabeling \(\Kstar\).\\
C & ontology-law research packet & not selected & Would add underived
\texttt{EqSrc}/allowedness semantics before the locality burden is decided.\\
D & protected human-gated stop & not triggered & A remains bounded,
non-promotional, and executable.\\
\bottomrule
\end{tabular}
\end{center}

\section{Selected future packet}

The selected identity is
\begin{center}
\small
\url{PKT-V22-P4T02-B2-KSTAR-LOCALIZATION-AND-BRIDGE-IRRELEVANCE-THEOREM-V1}.
\end{center}
Its registered packet type is
\texttt{source\_side\_irrelevance\_theorem}; its next role is
\texttt{refuter@0.2.0}.  It is selected but not executed.

The future Refuter must return exactly one of the following under explicit
hypotheses: a scoped \(\Kstar\) bridge-irrelevance theorem, a nontrivial
localization-compatible bridge-relevance witness, or the first precise
localization/factorization typing obstruction.  A negative theorem is limited
to the unchanged assignment and predeclared source-only bridge grammar.  It
must not be promoted to universal source-extension impossibility or global
theory rejection.

\section{Parent-child synthesis}

The Physicist--Mathematician and Physicist--Philosopher independently selected
Route A.  Both identified the compact/open mismatch and rejected a nominally
constructive Route B until a source-defined local codomain and map exist.  The
mathematical analysis contributed the extension-by-identity witness, the
first-jet representation hypotheses, and the quotient normality guard.  The
philosophical analysis contributed the strict distinction between scoped
irrelevance and a global no-go, and between mathematical orientation and
physical time orientation.  The parent retains all of these constraints.
There is no blocking conflict.

\section{Distance-to-GR and authority ledger}

No Distance-to-GR row changes.  Source primitives, \texttt{EqSrc}, retention,
generation, localization, response, \(M_{\rm src}\), \(g_{\rm eff}\), matter
coupling, Einstein equations, finite-variation robustness, benchmark status,
Gate status, and freeze status all remain \texttt{no\_delta}.  All six
inherited local freezes remain active.  D7 is not reevaluated, B2 is inactive,
and P4-T03 remains locked.

Validation, rendering, role labels, and checkpointing are process evidence
only.  They do not prove physics.  This transaction authorizes neither source
law or ontology adoption, nor proof promotion, publication, push, external
action, benchmark promotion, Gate review, or a completed derivation claim.

\end{document}
